% ==============================================================================
% Section 3: Merton (1973) Continuous Dividend Yield Formulation
% Target: Page 4 (Exactly 1 Page)
% ==============================================================================

\section{Merton (1973) Continuous Dividend Yield Formulation}
\label{sec:merton_dividend_model}

\subsection{Incorporation of Continuous Dividend Yield}

In standard corporate equity environments, high-dividend stocks pay regular cash disbursements. \citet{merton1973theory} extended the Black--Scholes framework to underlying assets paying a continuous dividend yield at constant rate $q \ge 0$. Under this model, an investor holding one unit of equity receives continuous dividend cash flow $q S_t \, \dd t$.

To maintain a self-financing delta-hedging portfolio $\Pi = V - \Delta S$, the dividend cash inflow generated by the underlying equity must be incorporated into the portfolio return dynamics:
\begin{equation}
    \dd \Pi = \dd V - \Delta \, \dd S - q \Delta S \, \dd t
    \label{eq:dpi_dividend}
\end{equation}
Setting $\Delta = \pdiff{V}{S}$ eliminates the stochastic diffusion term $\dd W_t$, yielding the deterministic portfolio evolution:
\begin{equation}
    \dd \Pi = \left( \pdiff{V}{t} + \frac{1}{2}\sigma^2 S^2 \frac{\partial^2 V}{\partial S^2} - q S \pdiff{V}{S} \right) \dd t
    \label{eq:dpi_det_dividend}
\end{equation}
Enforcing the no-arbitrage condition $\dd \Pi = r \Pi \, \dd t = r \left( V - S \pdiff{V}{S} \right) \dd t$ produces the \textbf{Merton Dividend-Adjusted Black--Scholes PDE}:
\begin{equation}
    \pdiff{V}{t} + (r - q) S \pdiff{V}{S} + \frac{1}{2}\sigma^2 S^2 \frac{\partial^2 V}{\partial S^2} - r V = 0
    \label{eq:merton_pde}
\end{equation}
The presence of dividend yield $q$ acts as an effective drag on forward asset growth, replacing the risk-free drift $r$ in the convection term with the net cost-of-carry $(r - q)$.

\subsection{Analytical Pricing Formulas for European Calls and Puts}

Solving Equation~\eqref{eq:merton_pde} subject to terminal boundary conditions yields the closed-form valuation formulas for European call price $C$ and European put price $P$ \citep{hull2018options}:
\begin{align}
    C(S_0, K, T, r, q, \sigma) &= S_0 e^{-qT} N(d_1) - K e^{-rT} N(d_2) \label{eq:merton_call} \\
    P(S_0, K, T, r, q, \sigma) &= K e^{-rT} N(-d_2) - S_0 e^{-qT} N(-d_1) \label{eq:merton_put}
\end{align}
where $N(\cdot)$ denotes the standard normal cumulative distribution function, and $d_1, d_2$ are given by:
\begin{equation}
    d_1 = \frac{\ln(S_0 / K) + \left( r - q + \frac{1}{2}\sigma^2 \right) T}{\sigma \sqrt{T}}, \qquad d_2 = d_1 - \sigma \sqrt{T}
    \label{eq:d1_d2_merton}
\end{equation}

\subsection{Dividend-Adjusted Put--Call Parity}

The foundational no-arbitrage relationship linking call and put premiums must be modified to account for the dividend yield $q$. Consider Portfolio A (Long 1 European Call + Cash $K e^{-rT}$) and Portfolio B (Long 1 European Put + $e^{-qT}$ shares of underlying stock).

At expiration date $T$, with dividends continuously reinvested in the underlying stock, Portfolio B grows to exactly 1 share of stock, giving terminal value $\max(K - S_T, 0) + S_T = \max(S_T, K)$. Portfolio A matures to $\max(S_T - K, 0) + K = \max(S_T, K)$. Because both portfolios generate identical terminal payoffs across all states of the world, they must command identical initial values today:
\begin{equation}
    C + K e^{-rT} = P + S_0 e^{-qT} \quad \iff \quad C - P = S_0 e^{-qT} - K e^{-rT}
    \label{eq:dividend_put_call_parity}
\end{equation}
Subtracting Equation~\eqref{eq:merton_put} directly from Equation~\eqref{eq:merton_call} confirms Equation~\eqref{eq:dividend_put_call_parity} analytically.

\clearpage
